\section{Length in small regions and repeated crossings}
\label{sec:11-slow-covering}

The fast-path estimate controls the distance a subwalk can travel in a
short time. We need the converse control: a subwalk cannot accumulate
much more than the $4/3$ power of its spatial span. This alone does not
bound all visits to a ball, since a walk may leave and return. After
proving the subwalk bound, we control the number of such returns by
joining disjoint crossings into cylinder polygons.

Throughout this section, $\alpha=3/4$, heights are measured in band units,
and $B_j$ denotes the total critical mass of strict port bridges of height
$j$, with $B_0=1$. We use
\[
 B_j\le C(1+j)^{\alpha-1},\qquad
 \sum_{j\le R}B_j\asymp (1+R)^\alpha,
\]
the irreducible-renewal probability law, the strip first-length moment,
and the full-box susceptibility bound established above. A bridge is
\emph{weak} if its endpoints attain its extreme heights; bounded extensions
and shortenings turn it into a strict port bridge. All such operations
below have bounded inverse multiplicity and alter critical weights by
bounded factors.

For a fixed lattice vertex $o$ and fixed $C_0<\infty$, write
\[
 \mu_H(\mathcal A)=
 \sum_{\substack{\gamma\text{ starts at }o\\
                  \operatorname{diam}(\gamma)\le C_0H}}
       \rho^{|\gamma|}\mathbf 1_{\mathcal A}(\gamma),\qquad H\ge2.
\]
Here \(|\gamma|\) counts edges, and vertex-subwalk lengths also count
edges. Port recodings change these counts and weights by bounded factors.
This is a finite, unnormalized measure. Its total mass is polynomial in
$H$ by the full-box bound. Simplicity also gives
$|\gamma|\le C H^2$ throughout its support. Constants may depend on $C_0$.
A bound described below as superpolynomial means that for every fixed
$A>0$ its mass is at most $C_AH^{-A}$, uniformly in $H\ge2$.

\Needspace{5\baselineskip}
\begin{theorem}[Uniform local length bound]\label{thm:11-local-length}
For every $\tau>0$ and $A>0$ there is $C_{\tau,A}<\infty$ such that,
for every $H\ge2$, outside a set of $\mu_H$-mass at most
$C_{\tau,A}H^{-A}$, every vertex subwalk $\sigma$ satisfies
\begin{equation}
 |\sigma|\le H^\tau(1+s(\sigma))^{4/3},
 \qquad s(\sigma)=\text{its vertical span in band units}.
 \label{eq:L-37}
\end{equation}
The same conclusion holds for each of the finitely many lattice
normal directions simultaneously.
\end{theorem}

The proof uses repeated cuts at turning extrema. A long path in a
small height range forces either a long chain of comparable pieces or
many separate pieces whose lengths are already large. Ordered renewal
probes make either configuration rare. We first handle two estimates
needed when a probe returns unusually close to its own start or a
piece has very small height.

\subsection{Returns inside one irreducible bridge}

\begin{lemma}[Small returns]\label{lem:11-small-return}
For every $c>0$ and $P>0$ there are $D<\infty$ and $C<\infty$ such that
if $I$ is sampled from the irreducible-renewal law and $J$ is its
height, then for every $r\ge2$,
\begin{equation}
 \Prob\!\left\{\begin{array}{l}
 J\in[r,2r),\ \text{$I$ returns to depth at most}
       \ r(1+\log r)^{-D}\\
 \text{after first reaching depth }cr
 \end{array}\right\}
 \le C r^{-\alpha}(1+\log r)^{-P}.
 \label{eq:T-36}
\end{equation}
The statement holds in either orientation of $I$.
\end{lemma}

\begin{proof}
Decreasing $c$ only enlarges the event, so assume that $c$ is a small
fixed constant. The middle-third extremal decomposition used for
\eqref{eq:F-33b}, now at the line of height $\lfloor cr\rfloor$, gives an
initial rising bridge $A_1$, a retreat $T$, and a final rising bridge
$A_2$. Place the initial endpoint at height zero. Let $y$ denote the height
of the last minimum after the first hit of this line, and let $z$
denote the height of a chosen maximum before that minimum, with a fixed
tie rule for its location along the path. Before recoding, their heights are $z$, $z-y$, and $J-y$.
For large $r$, the return assumption gives $y=o(r)$ while $z\ge cr+O(1)$
and $r\le J<2r$, so all three spans are comparable to $r$.
After recoding at the two turns, write their heights as $m_1,m_T,m_2$.
Choose dyads $ar$ and $br$ for $1+y$ and $1+J-z$, respectively. Thus
$1/r\le a,b\le C$, $a\le C(1+\log r)^{-D}$, and
\[
 |m_1-m_T|\le Car,\qquad |m_2-m_T|\le Cbr.
\]
The first irreducible of $A_1$ has height at least a constant times $ar$,
and the last of $A_2$ has height at least a constant times $br$. A
single-crossing seam below $y$, or above $z$, would otherwise be a seam
of $I$ itself: the bounded changes at the remote turn lie on the far side
of that seam. This would contradict irreducibility.

At each turn inspect the incident end strings to their first renewals
at depth at least $q$, where
$q\asymp r(1+\log r)^{-Q}$, with $Q$ fixed. Separate the required outer
irreducible at each unprobed end of $A_1,A_2$. For an outer bridge, type
$\mathrm M$ means that this outer jump and the string inspected from the
turn do not overlap; type $\mathrm Z$ means that they share a jump, which
must be the required outer jump. For $T$, inspect from both ends: type $\mathrm M$ means the two
stopped strings are disjoint, allowing a common cut, and type
$\mathrm Z$ means they share one irreducible. The three baseline masses are
\[
 r^\alpha(ar)^{-\alpha},\qquad r^\alpha,
 \qquad r^\alpha(br)^{-\alpha}.
\]
Each $\mathrm M$ piece leaves an independent residual bridge with freely
summed height at most $Cr$, supplying the factor $r^\alpha$ in its
baseline.

An inspected end is good, denoted $\mathrm G$, if its stopping renewal
lands below $c'r$, where $c'$ is sufficiently small depending on $c$.
Otherwise it is bad. A type $\mathrm Z$ retreat has two bad ends:
the two birth depths are below $q=o(r)$, so the jump they share spans
a fixed proportion of the retreat. If $f$ is the dyadic scale of one plus the birth
depth of its stopping jump, the prefix mass is at most $Cf^\alpha$ and
the jump has mass at most $Cr^{-\alpha}$. Consequently a bad end costs
$(f/r)^\alpha$ beyond baseline. There is only one jump variable when the two
inspections overlap.
For a type $\mathrm Z$ retreat its actual mass factor is
\[
 f_1^\alpha f_2^\alpha r^{-\alpha}
 =r^\alpha(f_1/r)^\alpha(f_2/r)^\alpha.
\]
For the first outer piece of type $\mathrm Z$, the shared jump is
macroscopic and its mass factor is instead
\[
 f^\alpha r^{-\alpha}
 =\bigl[r^\alpha(ar)^{-\alpha}\bigr](f/r)^\alpha a^\alpha.
\]
Thus the first outer piece has an extra factor $a^\alpha$ beyond its
baseline and bad-end factor; the second has the corresponding
$b^\alpha$. These equalities allocate the weight of one shared jump,
without replacing it by independent copies. The same additional
$a^\alpha$ is available when $A_1$ is $\mathrm M$ and its required first
jump has height at least $c'r$.

Use the logarithmic first-hit estimate \eqref{eq:T-35a} at each turn,
at a sufficiently small fraction of the smaller available depth: $q$
on a good end and $f$ on a bad end. At bounded depth use the trivial
bound instead. The test precedes every counted bad jump; after its
stopping renewal, the remaining prefix occurrences still cost at most
$Cf^\alpha$ by the renewal property. The two end prefixes at a turn
are distinct, and the recoding preserves their first-hit disjointness.
The two turn tests therefore give $r^{-\alpha}$ each, with the following
additional factors: $(r/q)^\alpha$ for two good ends, a constant for
one good and one bad end, and $(q/r)^\alpha$ for two bad ends. All other
logarithmic losses, including birth dyads and the choices of $a,b$, are
bounded by $(1+\log r)^{C_1}$, where $C_1$ does not depend on $D,Q$.

It remains to use the height constraints without conditioning the
probe laws. Put
\[
 K_r(j)=r^{-\alpha}B_j\mathbf 1_{\{0\le j\le Cr\}}.
\]
Uniformly under arbitrary translations, one kernel restricted to an
interval of width $w\ge1$ has mass at most $C(w/r)^\alpha$.
For $m=2,3$, H\"older's inequality gives
\[
 \sum_j\prod_{\nu=1}^m K_r(j+u_\nu)
 \le \prod_{\nu=1}^m\left(\sum_jK_r(j+u_\nu)^m\right)^{1/m}
 \le C r^{1-m}.
\]
Here $m(\alpha-1)>-1$, so the last estimate follows by summing the
power kernel. Summing the allowed height differences now shows that
two residual kernels matching within $Car$ cost $Ca$, and three
kernels with the two displayed height constraints cost $Cab$.
These bounds hold after every outer jump and every probe has been
fixed. If a residual height is itself at least $c''r$, its normalized
kernel is at most $C/r$, so restriction to width $Car$ costs $Ca$.

We apply these estimates in the following order. Fix all end strings
and distinct jump shapes, and first sum the residual bridge heights
with their translated restrictions. The kernel estimates just proved
are uniform in the fixed data. Then integrate every distinct separated
or shared jump once. Finally apply the first-hit estimates to the
remaining pre-jump end strings. Different turns use distinct such
strings. Their untested continuations contribute the uniform renewal
occurrence bounds already stated. This order proves simultaneous gains
without conditioning a probe law on a completed bridge height.

After the turn factors the baseline is $r^{-\alpha}(ab)^{-\alpha}$.
It remains to gain either a positive power of the small return scale
$a$ or a positive power of $q/r$. In the all-$\mathrm M$ case,
\[
 (ab)^{-\alpha}\,ab=a^{1-\alpha}b^{1-\alpha}\le C a^{1/4}.
\]
The other cases distribute the same kernel restrictions differently.
If $T$ and both outer pieces are $\mathrm M$, the triple-kernel bound
leaves $a^{1-\alpha}b^{1-\alpha}$. If $T,A_1$ are $\mathrm M$ and $A_2$
is $\mathrm Z$, use the two-kernel bound $a$ and its extra $b^\alpha$.
If $T$ is $\mathrm M$ and $A_1$ is $\mathrm Z$, use its extra $a^\alpha$,
restrict $T$ to width $Car$, and obtain $b^\alpha$ from either a single
restriction on $A_2$ or its $\mathrm Z$ improvement. In these three cases the remaining factor is bounded by
$C a^{\kappa}$, where $\kappa=\min(\alpha,1-\alpha)$, before
possible good--good losses.

If $T$ is $\mathrm Z$, neither turn has a good--good loss. Obtain
$b^\alpha$ from $A_2$ as above. When $A_1$ is $\mathrm Z$, its extra
$a^\alpha$ comes from the macroscopic shared-jump tail. When $A_1$
is $\mathrm M$ with a bad turn probe, fix the $\mathrm Z$ retreat's
height and use $|m_1-m_T|\le Car$ to restrict the free residual kernel
of $A_1$, obtaining $a^\alpha$. In both cases the tail and restriction factors cancel
$(ab)^{-\alpha}$; the remaining first-turn factor is explicitly
$(q/r)^\alpha$. When $A_1$ has a good turn
probe and a macroscopic first jump, the tail improvement and the
single residual restriction give $a^{2\alpha}$. In the remaining case
its first jump and good turn string both have height below $c'r$,
whereas $A_1$ has macroscopic height. Its residual height is therefore
macroscopic; the sharper single-kernel restriction gives $a$.
Thus, with $\kappa=\min(\alpha,1-\alpha)=1/4$, the total mass is at most
\[
 C r^{-\alpha}(1+\log r)^{C_1}
 \left[(1+\log r)^{-D\kappa+2Q\alpha}
              +(1+\log r)^{-Q\alpha}\right].
\]
Choose $Q$ first so that $Q\alpha>P+C_1$, and then $D$ so that
$D\kappa>P+C_1+2Q\alpha$. Bounded scales are absorbed by increasing $C$.
\end{proof}

The arbitrary logarithmic gain in this lemma will pay for early-birth
dyads at the many turns of a long chain. Without it, a power of
$\log H$ at every turn would overwhelm the later height-matching gain.

\subsection{Thin excursions}

\begin{lemma}[Thin excursions]\label{lem:11-thin-excursions}
For every $\lambda>0$ and $A>0$, there is $C_{\lambda,A}<\infty$ such that
\[
 \mu_H\!\left\{\begin{array}{l}
 \text{some subarc has vertical span at most }H^{\lambda/4}\\
 \text{and length greater than }H^\lambda
 \end{array}\right\}
 \le C_{\lambda,A}H^{-A},\qquad H\ge2.
\]
Subarcs may have vertex or horizontal-port endpoints, with bounded
endpoint changes in their lengths and spans.
\end{lemma}

\begin{proof}
If such an arc exists, simplicity and the lattice area bound force
horizontal range at least $cH^{3\lambda/4}$. One of the oblique lattice
normal projections consequently has range
$R\ge cH^{3\lambda/4}$. Between suitable extreme vertices of that
projection take a weak bridge, and extend its ends at bounded cost.
Its height is comparable to $R$ and its vertical range is still
$O(H^{\lambda/4})$. This is confinement in the direction tangent to
an oblique cut, rather than confinement in the oblique height.

Use the turning-extremum amplification in the proof of
Theorem~\ref{09_fast_paths:fast-bridges}, replacing its length cap
by this confinement. That argument turns a fixed power saving for
every free $\mathrm M$ residual bridge into arbitrary polynomial decay
by increasing a fixed iteration depth. We verify the only changed input, namely suppression of a
free $\mathrm M$ residual bridge. Let $V=O(H^{\lambda/4})=O(R^{1/3})$,
and suppose $R^{1-\kappa}\le s\le CR$, for a small fixed $\kappa>0$.
Residual heights $j\le R^{1-2\kappa}$ contribute at most
$CR^{\alpha(1-2\kappa)}$. For larger heights, the permitted vertical
endpoint difference gives at most $C(1+V)$ lateral endpoint offsets
at each height, since the tangent to the oblique cut has a nonzero
vertical component. If $b_j$ is the endpoint kernel, the variation
bound therefore gives
\[
 \sum_{R^{1-2\kappa}<j\le Cs}
       \sum_{\text{permitted offsets}}b_j
 \le C(1+V)\sum_{j\ge R^{1-2\kappa}}\|b_j\|_\infty
 \le C(1+V)B_{\lfloor R^{1-2\kappa}\rfloor}
 \le C R^{1/12+\kappa/2}.
\]
Both contributions are $O(s^\alpha R^{-a_0})$ for sufficiently small
$\kappa,a_0>0$. Every other factor in the turning-extrema amplification
is unchanged, so its arbitrarily large fixed iteration depth gives
mass $O_A(R^{-A})$. Sum over dyadic $R$, placements, and the two exterior
walk portions. Their cost is a fixed polynomial in $H$ by the full-box
bound; since $R$ is a positive power of $H$, it is absorbed by choosing
the preceding exponent large enough. Trimming port endpoints costs
only a bounded factor.
\end{proof}

\subsection{Turning chains and simultaneous integration}

We give the chain estimate with its probability bookkeeping, since the
number of turns will grow with $H$. Pieces in the following trees may
have vertex or horizontal-port endpoints. Their \emph{additive length}
counts a full edge as one and a half-edge as one half, so lengths add
exactly at every cut. This agrees with the vertex-subwalk convention
and with the visited-vertex count for a port-to-port path.
Fix a small $\epsilon>0$, set
$\beta=\epsilon/4$, and call a weak bridge of span $s$ \emph{long} when
its additive length exceeds $H^\epsilon(1+s)^{4/3}$. Fix a large integer
$p$ and a large constant $U$, to be chosen in that order. For every
eligible weak-bridge subarc of span $s\ge H^\beta$, form a random trial
tree as follows. At an expanded node choose a lattice line uniformly
from the middle third of its height interval. A single crossing gives
two forward children and an empty retreat, retained as a zero-length
child. Otherwise choose the last
minimum after the first hit of the line and a maximum before that
minimum, with fixed tie rules, obtaining the three chronological weak
bridges. Reverse or reflect coordinates when necessary. Expand each
child by the same rule, stopping when that child has span less than
the fixed threshold $s/U$, or at depth $L=\lceil\log H\rceil$.
Here $s$ remains the original root span throughout the tree. Choices
at distinct expanded nodes are fresh.

There are two kinds of witnesses. A \emph{spine witness} retains $L$
expanded nodes along one descent. A \emph{marked witness} is the minimal
subtree down to $p$ incomparable long nodes, which are retained as
marked leaves. In either case retain every unused child as a leaf, so
all expanded vertices have exactly three chronological children. If
$k$ is the number of expansions, then $k=L$ for a spine and $k\le pL$
for a marked witness. The ordered leaves form an alternating chain.

At a node with exactly one retained expanded child, call that child
\emph{continuing}; its two siblings are leaves. If the parent span is
$h$ and the three chronological child spans are $a,d,c$, then
$h=a-d+c$ and each child span is at most $h$. A small decrease from
$h$ to the continuing span forces its two sibling spans to be close,
as Figure~\ref{fig:11-sibling-heights} shows. The witness estimate will
use this constraint to obtain a small factor at many expansions.

\begin{figure}[ht]
\centering
\begin{tikzpicture}[x=.8cm,y=.65cm,font=\small,>=Stealth]
 \draw[gray!55,dashed] (0,0)--(6.8,0);
 \draw[gray!55,dashed] (2,3)--(6.8,3);
 \draw[gray!55,dashed] (2,4)--(6.8,4);
 \draw[very thick,blue!65!black,->] (0,0)--(2,3)
   node[midway,left,align=right] {$a$\\continuing child};
 \draw[thick,->] (2,3)--(4,1) node[midway,above right] {$d$};
 \draw[thick,->] (4,1)--(6,4) node[midway,right] {$c$};
 \draw[<->] (6.8,0)--(6.8,4) node[midway,right] {$h$};
 \draw[<->] (3,3)--(3,4) node[midway,left] {$h-a=c-d$};
 \node[below] at (0,0) {start};
 \node[above] at (6,4) {end};
 \node[below,align=center] at (4.5,.1) {two sibling leaves};
\end{tikzpicture}
\caption{A unary expansion with the first child continuing. The parent
and child heights obey $h=a-d+c$, so a small continuing-height drop
forces the two sibling leaves to have nearly equal heights. If the
retreat continues instead, $a,c\in[d,h]$ and $|a-c|\le h-d$.
Only heights are represented: the horizontal spacing and straight
segments are schematic.}
\label{fig:11-sibling-heights}
\end{figure}

For a fixed path $\gamma$, let $K_\gamma$ be the probability law of
all these independent trial choices, over every eligible root subarc
and every potential descendant node. This is a finite product: the
path has at most $CH^2$ steps and all trees have depth at most $L$.
For a trial outcome $\omega$, let $N^{\mathfrak t}_{r,k}(\gamma,\omega)$
count the witnesses of kind $\mathfrak t\in\{\mathrm{spine},\mathrm{mark}\}$
with $k$ expansions and root span in $[r,2r)$. Different root subarcs
and different retained subtrees are counted separately. The weighted
witness mass is therefore
\[
 \int N^{\mathfrak t}_{r,k}(\gamma,\omega)\,
       K_\gamma(d\omega)\,\mu_H(d\gamma).
\]
In particular, a bound on this integral bounds the mass of path-and-trial
pairs having at least one such witness.

\smallskip\noindent\emph{Why these witnesses suffice.}
Suppose that the thin-excursion exception of
Lemma~\ref{lem:11-thin-excursions}, with $\lambda=\epsilon$, is absent,
and that a trial outcome contains neither kind of witness, for any
eligible root. If
\[
 pU^{-4/3}<1/4,
\]
then, for all sufficiently large $H$, every weak-bridge subarc $\sigma$
of span $s$ satisfies
\begin{equation}
 |\sigma|\le H^{2\epsilon}(1+s)^{4/3}.
 \label{eq:11-no-witness-length}
\end{equation}
This implication is deterministic; it holds for every outcome of the
trial choices.

Prove it by induction on three times the band span, an integer for
both vertex and horizontal-port endpoints. Below $H^\beta$ the thin
bound gives length at most $H^\epsilon$. For a larger root span $s$,
take its trial tree and prune at each non-long piece. At any depth
fewer than $p$ long pieces remain, since they are incomparable and
would otherwise give a marked witness. Thus at most $1+3pL$ non-long
pieces stop the pruning. Every long terminal piece has span below
$s/U$: a terminal at depth $L$ would give a spine. These long terminals
are also incomparable, so there are fewer than $p$ of them. Their
spans are strictly smaller than $s$, and the induction hypothesis
and exact length additivity give
\[
 |\sigma|\le (1+3pL)H^\epsilon(1+s)^{4/3}
       +pH^{2\epsilon}(1+s/U)^{4/3}.
\]
After division by the right-hand side of
\eqref{eq:11-no-witness-length}, the coefficient is
\[
 (1+3pL)H^{-\epsilon}
 +p\left(\frac{1+s/U}{1+s}\right)^{4/3}<1
\]
for large $H$, uniformly over $s\ge H^\beta$. This proves the
induction. It remains to show that the weighted mass of the two
witness classes is small.

\begin{lemma}[Turning-chain witness bound]\label{lem:11-chain}
There are absolute $b>0$ and $C<\infty$ such that, for fixed
$\epsilon,p,U$ and a root-span dyad $s\in[r,2r)$ with
$H^\beta/2\le r\le C_0H$, the weighted witness mass just defined is at most
\begin{equation}
 H^{C+o_{p,U,\epsilon}(1)}H^{-\epsilon m_*}
 \bigl[C'(1+\log\log H)^{C'}\bigr]^k(k+1)^{-bk}.
 \label{eq:L-39}
\end{equation}
There $m_*=0$ for a spine and $m_*=p$ for a
marked witness. The exponent
$C$ is independent of $p,U,\epsilon$; the constant $C'$ may depend on
them. The $o(1)$ is uniform in the indicated $r,k$, and denotes a
quantity tending to zero for those fixed parameters. The same bound,
with $m_*=0$ and $k$ equal to the number of pieces, holds for an
alternating chain of $1\le k\le pL$ pieces of nonincreasing spans in
$[r,2r)$, where $L=\lceil\log H\rceil$ and $p$ remains fixed.
For this last assertion the mass is the analogous sum over all
rooted chain records, without trial probabilities.
\end{lemma}

\begin{proof}
\emph{Cuts and bridge weights.}
We first record the cuts and relative bridge shapes. The calculation
then separates end tests from the remaining height variables. The
height restrictions are summed by bounds uniform in the end data, so
the tests can be integrated with their original renewal laws.

Every outer child of an expansion has span at least one third
of its parent's span. Choose $0<c_0\le1/(12U)$. Every expanded parent
has span at least $s/U\ge r/U$, so all leaves have spans between
$c_0r$ and $2r$, except possibly retreat leaves.
There are only $O_{p,U}(1)$ small retreat leaves, and they are isolated
between macroscopic leaves. To see the bound, the minimal witness has
$O(p)$ branching or terminal expansions and otherwise consists of
$O(p)$ unary chains. If a parent has height $h$ and a small retreat
has height $d<c_0r$, the middle-third cut gives both forward heights
at most $2h/3+d\le3h/4$. The retreat itself has height below $s/U$
and cannot continue.
Along a continuing chain heights never increase and stay at least
$s/U$ until the last expansion; there can consequently be only
$O_U(1)$ such reductions on each chain.
The first and last leaves of an expanded subtree descend through
outer children and are macroscopic, so small retreats are isolated
in the chronological leaf order.

Bin each small span plus one dyadically with scale $r_i$, and put
$r_i=r$ on macroscopic leaves. Fix the complete leaf pieces and all
cut records. At the parent of a small retreat, the chosen middle-third
line must belong to that retreat's height interval (one specified line
when it is empty); there are $O(r_i)$ possible choices out of at least
$cr$. Each small leaf therefore contributes $Cr_i/r$ to the probability
of these witness data. All other line-choice probabilities sum to at
most one. This uses only choices recorded by the witness, without
specifying any descendants outside it.

At a turn whose incident spans are both larger than $Cr^{1/20}$,
recode its incident leaves to strict bridges from neighboring ports,
preserving disjoint first-hit portions to distance
$c_A\min(r_i,r_{i+1})$. Here $c_A$ is small depending on $c_0$.
For completeness, at a peak both lower neighbors of its triangle occur
on the path. One arm uses the central top port. Launch the other from
the adjacent top port through the neighboring downward triangle to
its original lower neighbor. If that neighboring triangle already
occurs on this arm, it must occur immediately next to that lower
neighbor, so shorten the initial segment instead. An internal visit
would have to use the same lower neighbor. The other arm cannot use
this triangle in the retained first-hit portion for the same reason.
Reflect this operation at minima. Remote-end changes lie beyond the
retained first-hit portions. Where no turn test is required, extend
to bounding port lines; if an extreme triangle does not border that
line, one extra triangle nearer the line is unused by extremality.
Empty leaves have bounded bookkeeping. This gives a finite record
per cut and changes lengths and spans by $O(1)$.

After fixing those records, normalized bridge shapes have their
translations fixed successively by incidence. There is no additional
sum over horizontal positions at each turn. A turn test involves only
two relative shapes launched at fixed neighboring ports; changing
absolute translations does not change the test. The root location,
the two exterior portions, and endpoints cost $O(H^C)$ altogether by
the full-box bound, with $C$ independent of the number of leaves.

Baseline leaf $i$ by $r_i^\alpha$ and initially seek
$\min(r_i,r_{i+1})^{-\alpha}$ at every turn. Because small leaves are
isolated, multiplying the baselines, these turn factors, and the
placement probabilities gives at most
\[
 C^k r^\alpha\prod_{i\text{ small}}(r_i/r)^{1-\alpha}.
\]
If $r_i\le Cr^{1/20}$, omit both adjacent turn tests. Restoring their
nominal factors multiplies this expression by at most $Cr_i^{2\alpha}$.
The resulting contribution from that leaf is bounded, since
\[
 (r_i/r)^{1-\alpha}r_i^{2\alpha}
 \le C r^{-(1-\alpha)+(1+\alpha)/20}.
\]
No probes are needed at its own ends. Call turns touching small or
marked leaves \emph{rough}; there are $O_{p,U}(1)$ of these. All other
turns are \emph{ordinary}.

\smallskip\noindent\emph{End probes at ordinary turns.}
The factors at turns will cancel all but one of the bridge mass scales.
Their dependence on $k$ must remain explicit, because $k$ can be of
order $\log H$.

For a non-tiny leaf inspect its two strings of irreducibles through
the first renewals reaching
\[
 q_i=\left\lfloor c_qr_i(k+1)^{-\delta}\right\rfloor,
\]
with a sufficiently small fixed $\delta>0$. The coefficient $c_q$ is
small, and common on macroscopic leaves. For large $H$ these are
positive. The strings are of type $\mathrm M$ when disjoint, allowing
a common terminal cut, and then leave a free middle bridge. Otherwise
they share exactly one irreducible and have type $\mathrm Z$. Call an
end good when its stop lands below $c'r_i$, and bad otherwise, with
$c_q\ll c'\ll c_Ac_0$. In type $\mathrm Z$ both ends are bad.
A bad string is a prefix with endpoint depth below $q_i$, followed by
a jump of height comparable to $r_i$. For a birth dyad $f$, its prefix
cost is $Cf^\alpha$ and its jump cost is $Cr_i^{-\alpha}$. On an
$\mathrm M$ leaf the residual middle supplies baseline $r_i^\alpha$;
on a $\mathrm Z$ leaf its single shared jump supplies
$r_i^\alpha(r_i^{-\alpha})^2=r_i^{-\alpha}$. This equality is only
bookkeeping: every distinct shared jump is integrated once.
For example, among $n_0=m_0+z_0$ macroscopic pieces with $b_0$ bad
ends there are $d_{\mathrm{bad}}=b_0-z_0$ distinct bad jumps, and
\[
 r^{\alpha n_0}r^{-\alpha b_0}
       =r^{\alpha m_0}r^{-\alpha d_{\mathrm{bad}}}.
\]
The right-hand side records the genuine independent middle and jump
variables; the left-hand side assigns their costs to leaves and turns.

At an ordinary turn the desired estimates, beyond these leaf
baselines, are
\begin{equation}
\begin{array}{c|c}
 \text{incident ends}&\text{turn factor}\\ \hline
 \mathrm{GG}&Cr^{-\alpha}(k+1)^{\delta\alpha}\\
 \text{one good, one bad}&Cr^{-\alpha}\\
 \text{two bad}&Cr^{-\alpha}(k+1)^{-\delta\alpha}.
\end{array}
\label{eq:L-38}
\end{equation}
The total additional loss over all ordinary turns will be at most
$C^k(1+\log\log H)^{Ck}$.

First fix the birth dyads. Apply the ordered test $E$ from
\eqref{eq:T-34} at a small fraction of the minimum of the available
depths, using $q_i$ at a good end and $f$ at a bad end. At bounded
minimum use the trivial test. Its stopping renewals occur before the
first-hit portions terminate: on a bad end they precede the specified
birth, and a good stop is below $c'r$. Renewal cuts prevent later
visits below the tested level. Moreover two disjoint crosscuts from
the neighboring launches to the common first-hit line have the
required left--right extremal order at each lower level; a point past
the opposite extremum would lie on the wrong exterior side of that
crosscut. Hence $E$ is a necessary condition. After these stops the
remaining birth counts are at most $Cf^\alpha$, including the
possibility that the birth occurs at the stop itself. The three costs
for fixed dyads are respectively
\[
 Cq^{-\alpha},\qquad
 Cr^{-\alpha}f^\alpha\min(f,q)^{-\alpha},\qquad
 Cr^{-2\alpha}f^\alpha g^\alpha\min(f,g)^{-\alpha}.
\]
Since $f,g\le Cq$ these give the table; the last expression is more
precisely $Cr^{-\alpha}(\max(f,g)/r)^\alpha$.

\smallskip\noindent\emph{Summing birth depths.}
We must sum the birth dyads without a power of $\log H$ at every turn.
Set
\[
 t=\lfloor r(1+\log H)^{-D_*}\rfloor,\qquad
 T_0=r(1+\log H)^{-10},
\]
where $D_*>12$ is a sufficiently large fixed constant. Tag an unmarked
macroscopic bad jump if, from either of its relevant ends, it returns
to depth at most $2t$ after reaching $c''r$, where $c''>0$ is small.
By Lemma~\ref{lem:11-small-return}, choosing $D_*$ larger if necessary,
the mass of tagged jumps is
\[
 Cr^{-\alpha}(1+\log H)^{-5}.
\]
The constants are uniform for $H^\beta/2\le r\le C_0H$, since
$\log r\asymp_\beta\log H$. The union over two ends costs only a
constant, and no designated-jump test below will be assigned to a
tagged jump.

If both bad birth depths are at most $T_0$, sum their dyads using the
improved maximum factor. Summing over the larger dyad first gives
\[
 \sum_{f,g\le CT_0}\left(\frac{\max(f,g)}r\right)^\alpha
 \le C(1+\log H)(T_0/r)^\alpha
 \le C(k+1)^{-\delta\alpha},
\]
for small $\delta$, since $k\le pL$. Birth dyads between $t/L_2$ and
$Cq$ have only $O(1+\log\log H)$ possible values. Here $L_2$ is the
fixed paired-test constant in \eqref{eq:T-34}, unrelated to the trial
depth $L$. The remaining situation has a birth below $t/L_2$ opposite
a good end or a bad birth larger than $T_0$. If its jump is tagged,
sum this early birth dyad directly. A distinct jump has at most two
ends; its factor $(1+\log H)^{-5}$ pays both possible lost dyad sums,
each of size $O(1+\log H)$.

For an untagged jump in that remaining situation, use
$F(t,r',0)$ from \eqref{eq:T-34}, designating the early-birth end;
$r'$ is the jump's height dyad, with only a bounded number of choices
comparable to $r$. The opposite end supplies an independent worker
string. Its test stop precedes any specified bad birth there, since
that birth is above $T_0\gg t$. On the designated arc, every visit
below $t$ up through the jump occurs before the common first-hit line.
Otherwise, as the designated birth is below $t/L_2$, the jump itself
would reach $c''r$ and subsequently return below $2t$, forcing its tag.
The crosscut argument thus gives the required $E(t)$ event. If the
same jump is designated from both ends, use the single paired bound
$F_{12}(t,t,r')$ instead. It integrates that one jump, both early
prefixes, and both opposite workers' tests. In either case the cost
is $Cr^{-\alpha}$ for the distinct jump and all its designated
prefixes and tests. A worker with a later bad birth still contributes
$Cg^\alpha$ after its test, so a two-bad turn retains its factor
$(g/r)^\alpha\le C(k+1)^{-\delta\alpha}$. The single-test birth
condition holds because $L_2$ is at least the single-test ratio in
\eqref{eq:T-34}.

\smallskip\noindent\emph{Simultaneous integration.}
We will combine the turn gains with restrictions on residual middle
heights. The needed restriction can be stated now. For an unmarked
macroscopic $\mathrm M$ leaf, let $z$ be the combined height of its two
end strings and $j$ its middle height. Its total height is $z+j$.
The normalized middle kernel
\[
 K_r(j)=r^{-\alpha}B_j\mathbf1_{\{0\le j\le Cr\}}
\]
obeys, for every interval $I$ of width $w\ge1$ and every shift $z$,
\begin{equation}
 \sum_{j:\,z+j\in I}K_r(j)\le C(w/r)^\alpha.
 \label{eq:11-middle-window}
\end{equation}
Indeed $B_j\le C(1+j)^{\alpha-1}$, so the sum over any $w+1$
consecutive heights is at most $Cw^\alpha$. Since $k\le pL$ and
$r\ge H^\beta/2$, requiring
two leaf heights to differ by $O(1+r/\sqrt{k})$ costs
$O((k+1)^{-\alpha/2})$ if one leaf is of type $\mathrm M$: fix the
partner's height, sum that middle in its translated window, and then
sum the partner's middle freely if it has one. The bound is uniform
in the partner's end strings and jumps, leaving those variables
available for their later probe estimates. For disjoint leaf pairs
these middle integrations can be performed together.
The construction of sufficiently many such pairs is given below.

We now integrate the turn tests while retaining these uniform
middle-height gains. Work with nonnegative critical-weight sums and
fix only discrete records: the chain shape, scales, piece types,
good/bad labels, tag labels, and assignments of turn tests. Do not
condition any bridge law on its total height or on the completed
chain. For an $\mathrm M$ leaf its atomic variables are its two
oriented end prefixes, any distinct bad stopping jumps, and its
middle bridge. For a $\mathrm Z$ leaf they are the two pre-jump
prefixes and one shared jump. Concatenation identifies a tuple of
such variables with at most a constant number of original records.
Dropping unused self-avoidance and incidence conditions therefore
bounds the original sum by a product of their critical measures.
The restrictions retained in this product are local turn tests,
individual jump restrictions, and the residual-height constraints
just described.

Each end prefix occurs at exactly one turn, or at a free chain end.
Assign every distinct bad jump to exactly one of three blocks: its
individual tail bound (with its tag restriction if present), a
single $F$ block, or a paired $F_{12}$ block. A designated block
contains the designated pre-jump prefix or prefixes and the opposite
worker prefixes through their test stops. An ordinary $E$ block
contains its two prefixes only through their test stops. These sets
of variables are disjoint. In particular a worker at a designated
turn has no early designated birth there: it is either good or its
birth is above $T_0$. The tested worker portion therefore excludes
its own separately charged bad jump. If a shared jump has just one
assigned $F$ test, the test at its other end uses only the prefix
before that jump; it belongs to a different block and never uses
the shared jump a second time.

First sum the residual middles with the prescribed height
restrictions, using \eqref{eq:11-middle-window} for one middle in
each disjoint pair and the unrestricted bound for the others.
These bounds are uniform in all the end variables. Leave every designated
pre-jump prefix intact inside its assigned $F$ or $F_{12}$ block.
For an ordinary $E$ prefix or a worker prefix, next sum only its
untested continuation. Given its test stopping history,
the mass of subsequent renewal occurrences in a birth dyad $g$ is
at most $Cg^\alpha$, uniformly in the stopping location and overshoot:
an overshoot beyond the permitted endpoint contributes zero, and
otherwise the remaining range has length at most $Cg$. A good
completion is a probability at most one. Length insertions are
handled by the same conditional argument described below. Replacing
these continuation sums by their uniform bounds removes them without
conditioning the preceding tests. Now integrate each disjoint $E$,
$F$, or $F_{12}$ block using \eqref{eq:T-34}; integrate every remaining
jump using its individual tail bound. Tonelli's theorem justifies this
order. Relative launches in each block are fixed neighboring ports,
so absolute translations depending on variables outside the block
introduce no parameter into its estimate. In a paired block this is
exactly the two-end convention in the definition of $F_{12}$.
Finally sum the discrete labels and dyads as above. Thus no estimate
has been applied to a renewal law conditioned on an eventual height,
a tag, another turn's survival, or full-chain avoidance. This proves
\eqref{eq:L-38} with the stated total loss.

\smallskip\noindent\emph{Long marked leaves.}
The ordinary-turn estimates control the unmarked chain. A marked leaf
must additionally pay $H^{-\epsilon+o(1)}$ for exceeding its length
threshold. Only a fixed number of turns touch such leaves, so the
logarithmic first-hit estimate suffices there.

We next include rough turns and the marked length factors. Uniformly
for $1\le u\le C H$, the first-length moment gives
\[
 \sum_{j\le Cu}\sum_{B:\,\operatorname{ht}(B)=j}
       \rho^{|B|}(1+|B|)
 \le u^{\alpha+4/3}H^{o(1)}.
\]
The analogous summed moment for an irreducible of height at most $Cu$
is at most $u^{4/3-\alpha}H^{o(1)}$. Indeed, in a height dyad $b$,
prepend and append bridges of heights in $[b,2b]$, whose summed
masses each are at least $cb^\alpha$. The resulting bridge has height
$O(b)$ and has only boundedly many irreducibles of height at least $b$.
Its length moment bounds the marked jump's length moment, giving
$b^{4/3-\alpha}H^{o(1)}$ after division by $b^{2\alpha}$.
Sum the dyads, using $4/3>\alpha$. These moment statements mean that
for every fixed $\eta>0$ the displayed factors $H^{o(1)}$ are bounded
by $C_\eta H^\eta$, uniformly over all indicated scales.

Expand the product of length-plus-one factors on marked leaves by
assigning each insertion to its middle, a bad jump, a good end string,
or a pre-bad-jump prefix. There are only boundedly many categories
per marked leaf; for an end string sum over its individual jumps.
A middle or tiny-piece insertion costs at most
$r_i^{4/3}H^{o(1)}$ beyond its unweighted baseline. A bad-jump insertion
has the same relative cost by the irreducible moment. No tag saving
or designated $F$ test uses a marked leaf's own jump.

For an insertion in an end string, let $d$ be the dyadic scale of one
plus the inserted jump's birth depth. Up to that birth the occurrence
mass is at most $Cd^\alpha$. The jump's length moment, after extracting
$r_i^{4/3}H^{o(1)}$, costs $Cr_i^{-\alpha}$. Thus this insertion gives
an additional factor $(d/r_i)^\alpha$. Any subsequent required bad
birth of dyad $f$ costs at most $Cf^\alpha$ by restarting after the
inserted jump, with its bad-jump price already separately budgeted;
completion to a good stop has probability at most one. This includes
the terminal stopping jump of a good string. Additive constants in
the length weight use the unweighted estimate.

At a non-omitted rough turn apply the first-hit bound
\eqref{eq:T-35a} at a small fraction of the smaller available depth.
Use $d$ on an end with an insertion in its prefix, $f$ on another bad
end, and $q_i$ on another good end. The test precedes the counted
births and lies within the preserved disjointness range. Subsequent
occurrence counts have the same conditional bounds as before. A
side of available depth $v$ has, after extraction of its baseline
and any length price, cost at most
$H^{o(1)}(v/r_i)^\alpha$. For a good side the factor is supplied at
polylogarithmic cost, since $q_i/r_i$ is a negative power of $k+1$.
For two incident sides the test and these factors give
\[
 H^{o(1)}\frac{v_1^\alpha v_2^\alpha}
                   {r_i^\alpha r_{i+1}^\alpha}
             \min(v_1,v_2)^{-\alpha}
 \le H^{o(1)}\min(r_i,r_{i+1})^{-\alpha},
\]
because $v_j\le Cr_j$. Free ends and omitted turns cost at most
$H^{o(1)}$. There are only $O_{p,U}(1)$ rough turns and small leaves.
Their scale choices cost $(1+\log H)^{O_{p,U}(1)}$, so all these losses
remain $H^{o_{p,U}(1)}$.

This rough calculation is compatible with the preceding block
integration. At a rough turn an unmarked neighboring leaf supplies
only its good prefix or its prefix before a bad jump; no marked
insertion or rough test consumes that unmarked bad jump. Stop the
rough tests before the inserted jumps, then use the conditional
moment and renewal bounds for the continuations. Distinct marked
leaves have distinct inserted jumps. Their length prices hence
multiply. On a long marked leaf,
\[
 \mathbf 1_{\{|B|>H^\epsilon(1+\operatorname{ht}(B))^{4/3}\}}
 \le \frac{1+|B|}{H^\epsilon(1+\operatorname{ht}(B))^{4/3}}.
\]
Its scale is comparable to $1+\operatorname{ht}(B)$, so division by
this threshold gives $H^{-\epsilon+o(1)}$ per mark.

\smallskip\noindent\emph{The gain from many leaves.}
We have obtained the turn baselines and the cost of each marked
length. It remains to obtain the factor $(k+1)^{-bk}$ from the many
unmarked macroscopic leaves. We now construct the close-height pairs
to which \eqref{eq:11-middle-window} applies.

There are disjoint sibling pairs covering all but
$O_{p,U}(\sqrt{k}+1)$ leaves for which the two leaf heights differ by
$O(1+r/\sqrt{k})$. Indeed all but $O(p)$ expansions lie on unary
chains. The total drop of continuing heights along each such chain
is at most $2r$, so at most $O(\sqrt{k})$ expansions on that chain
have a drop larger than $r/\sqrt{k}$. At every other unary expansion
the two non-continuing children are
leaves. Write $h$ for the parent height, $a,c$ for the forward child
heights, and $d$ for the retreat height. The signed-height identity is
$h=a-d+c$, and all three child heights are at most $h$.
If the first child continues, the drop $h-a=c-d$ is exactly the
difference of the other two heights; the last-child case is symmetric.
If the retreat continues, then $a+c=h+d$, with $a,c\le h$, so
$a,c\in[d,h]$ and $|a-c|\le h-d$. Consequently a continuing-height
drop at most $r/\sqrt{k}$ always supplies a sibling pair whose heights
differ by at most that amount, up to the bounded port recodings. Remove
the boundedly many pairs touching marked or small leaves. Specifying
which of these pairs meet the constraint and their types costs at
most exponentially many choices in $k$.

In each retained pair with at least one $\mathrm M$ leaf, integrate
one such middle kernel with the height constraint, holding the
partner's height fixed. If the partner also has a middle, sum that
middle without the constraint. The pairs
are disjoint, so these integrations are simultaneous and uniform in
all probe variables. These are the residual integrations used in
the block argument above. There are $N=2k+1$ leaves. If at least $N/10$ are $\mathrm M$,
then, after removing the $O_{p,U}(\sqrt{k}+1)$ exceptions, at least
$N/40$ retained pairs contain an $\mathrm M$ leaf for large $k$.
Their total gain is at most $(k+1)^{-\alpha N/80}$. The at most $N$
good--good turns cost at most $(k+1)^{\delta\alpha N}$, so a fixed
choice $0<\delta<1/160$ leaves a negative multiple of $k\log(k+1)$.
If fewer than $N/10$ leaves are $\mathrm M$, at least
$4N/5-1-O_{p,U}(1)$ adjacencies join two $\mathrm Z$ leaves,
allowing for the bounded number of small leaves and rough exceptions. At most $N/10$ adjacencies can be good--good.
For large $k$ the excess of bad--bad over good--good ordinary turns
is therefore at least $N/2$, giving the factor
$(k+1)^{-\delta\alpha N/2}$.
Thus in either case the total factor
is at most $(k+1)^{-bk}$ for an absolute $b>0$, apart from constants
to the power $k$ and the stated logarithmic losses. Bounded $k$ are
absorbed by changing the constants.

Ordered ternary tree shapes, choices of the continuing children,
piece labels, and the finite recoding records have at most
exponential multiplicity in $k$. Choices of the fixed number of
marks contribute $H^{o_p(1)}$. Small-scale dyads do so as well. The
factor $r^\alpha$, root placements and exterior paths are all covered
by a fixed polynomial $H^C$ independent of $p,U$. This proves the bound in \eqref{eq:L-39}.

For the final assertion of the lemma, start instead with the given
alternating chain of $k\le pL$ comparable, nonincreasing spans. This
range preserves the probe and dyad estimates above. There are no
small leaves or placement weights. Pair successive pieces, leaving
at most one unpaired. Since the total height drop is at most $2r$,
all but $O(\sqrt{k})$ pairs have difference at most $Cr/\sqrt{k}$.
The same recoding, tests, variable partition, and middle-kernel
integrations apply, and give the stated estimate.
\end{proof}

\subsection{Slow-path amplification}

\begin{proof}[Proof of Theorem~\ref{thm:11-local-length}]
Choose $\epsilon>0$ much smaller than $\tau$. In the notation of
Lemma~\ref{lem:11-chain}, when $k=L$ the logarithm of the last two
factors of \eqref{eq:L-39} is
\[
 -bL\log L+O\bigl(L\log(2+\log\log H)\bigr),
\]
which is smaller than $-A\log H$ for every fixed $A$ and all
sufficiently large $H$. Thus spine witnesses have superpolynomially
small mass. For arbitrary $k\le pL$, let
$V_H=C'(1+\log\log H)^{C'}$. The maximum of
$k\log V_H-bk\log(k+1)$ is $O(V_H^{1/b})=o(\log H)$.
Consequently the sum over $k$ costs $H^{o(1)}$. Choose $p$ sufficiently
large that the factor $H^{-p\epsilon}$ absorbs the fixed polynomial
and the desired failure power. Next choose $U$ large enough that
$pU^{-4/3}<1/4$. Constants depending on these now fixed parameters
do not affect either conclusion. Summing the $O(\log H)$ root dyads
makes all unwanted witnesses negligible at any prescribed power.

Sum the witness bounds over their sizes and use the kernel
$K_\gamma$ defined above. Since each witness count dominates its
existence indicator, the resulting bound controls the
$\mu_H(d\gamma)K_\gamma(d\omega)$-mass of pairs containing any unwanted
witness. The root-subarc choices are already included in that count.
Add the path exceptional set from
Lemma~\ref{lem:11-thin-excursions}, with parameter $\lambda=\epsilon$.

Every remaining path-and-trial configuration satisfies the
deterministic implication \eqref{eq:11-no-witness-length}.
That weak-bridge statement is a property of $\gamma$ alone.
If a path violates it and is not in the thin exceptional set, every
trial configuration over that path must contain an unwanted witness.
Therefore, writing $\mathcal B$ for the bad paths and $\mathcal W$
for the unwanted configurations,
\[
 \mu_H(\mathcal B)
 =\int_{\mathcal B}K_\gamma(\Omega_\gamma)\,d\mu_H(\gamma)
 \le (\mu_H K)(\mathcal W)
       +\mu_H(\text{thin-excursion exception}).
\]
There is no loss from projecting away the trial variables.

To pass to arbitrary subwalks, split at a lowest vertex and read both
halves away from that vertex. Successively cut each half at the last
alternating extremum of the remaining tail. This gives weak bridges
of nonincreasing spans. Outside another superpolynomial exceptional
set, there are fewer than $\lceil\log H\rceil$ consecutive pieces in
any one span dyad above $H^\beta/2$: otherwise apply the final part
of Lemma~\ref{lem:11-chain} to $k=L=\lceil\log H\rceil$ of those
pieces. Root locations and
exterior pieces again have only fixed polynomial cost. Summing the
weak-bridge length bound over decreasing dyads gives at most
$CH^{2\epsilon}\log H(1+s)^{4/3}$ per half. Once the remaining span is
below $H^\beta$, the whole remaining tail is thin and has length at
most $H^\epsilon$. Choose $\epsilon$ so that these bounds fit
$H^\tau(1+s)^{4/3}$ for large $H$. Enlarging the exceptional-set
constant handles bounded $H$. A finite union over lattice symmetries
proves the simultaneous directional statement.
\end{proof}

\subsection{Many crossings}

We use the all-count cylinder estimate stated among the analytic
inputs. For clarity, its precise physical form is the following: on
the balanced two-marker cylinder with $N=2m$ bands, the first $m$ band
parameters are $2\lambda$ and the next $m$ are $\lambda$. If $w_N(a)$
is the critical mass of families of exactly $a$ mutually disjoint
polygons separating the markers, then
\begin{equation}
 w_N(a)\le
 \exp\!\left(C\log N-\frac{ca^2}{\log N}\right)
 \qquad(a\ge0)
 \label{eq:L-40}
\end{equation}
for every sufficiently large even $N$, with one pair of constants
independent of $N$ and $a$. The analytic input proves, before sector
truncation, $Z_N(2\cosh v)\le\exp(C\log N(1+v^2))$ for every
$v\ge0$; positivity and optimization in $v$ give
\eqref{eq:L-40}. The quantifier over all $a$ is needed below, since
the number of crossings grows as a power of the ambient scale.

\begin{lemma}[Many bridges in one slab]\label{lem:11-many-bridges}
Fix $D<\infty$. Consider ordered tuples of $a$ mutually avoiding
strict bridges across the same $j\ge1$ bands, all contained in one
abscissa slab of width $Dj$. Prescribe the distinct starting ports,
the distinct ending ports, and their matching. The critical mass of
these tuples is at most
\begin{equation}
 \exp\!\left(C_D\log(2+j)-\frac{c_Da^2}{\log(2+j)}\right).
 \label{eq:L-41}
\end{equation}
\end{lemma}

\begin{proof}
Put $d_0=\sqrt3/2$ and use triangular coordinates with generators
$(1,0)$ and $(1/2,d_0)$. Choose a fixed integer $b=b(D)$ sufficiently
large and put $m=bj$. In the physical cylinder the successive half-period
translations are
\[
 m(-1/2,d_0),\qquad m(1/2,d_0).
\]
Thus the total period is $(0,2md_0)$ and the marker abscissas differ
by $m/2$. For $b$ sufficiently large, an integer horizontal translation
places the entire width-$Dj$ slab strictly between those abscissas,
with a fixed positive margin. Align its bottom with a marker band level.

Label the fixed lower ports $s_1,\ldots,s_a$ and the fixed upper
ports $t_1,\ldots,t_a$, and write $\pi$ for the fixed matching:
strand $i$ joins $s_i$ to $t_{\pi(i)}$. Choose independently $2b$
tuples from the ensemble. Put the first in bands $0,\ldots,j$ and
reflect the second across the line at height $jd_0$, traversing it
upward. Reflection across any horizontal band boundary is a lattice
symmetry. The second strip takes $t_{\pi(i)}$ to the translate of
$s_i$ at height $2jd_0$. Repeat this construction in successive double
strips. Each double strip therefore induces the identity on the lower
port labels. After all $2b$ strips, the vertical quotient closes each
strand at its own initial port.

The strip interiors are disjoint. Strict bridges meet their boundary
cuts only at the fixed ports, and those ports are distinct within each
tuple. Hence the resulting family contains exactly $a$ disjoint simple
polygons. Each lift travels through one full vertical period, so each
polygon is essential. Since the whole family stays between the marker
abscissas, each polygon separates the two markers.

Cutting this family at the fixed $j$-spaced levels recovers every input
tuple. The strip index specifies its reflection, and the prescribed
ports recover the strand labels and matching. Recovery is therefore
injective even though the polygon family itself is unlabeled and
unoriented. Ports carry no vertex weight, so the critical weights
multiply exactly at every concatenation. If $M$ is the input tuple
mass, positivity and \eqref{eq:L-40} give
\[
 M^{2b}\le w_{2bj}(a)
 \le\exp\!\left(C\log(2bj)-\frac{ca^2}{\log(2bj)}\right).
\]
Take the $2b$th root and use
$\log(2bj)\asymp_b\log(2+j)$. For bounded $j$, the fixed-width slab
contains only boundedly many distinct ports, so $a$ is bounded;
the finite bridge masses cover these cases by changing the constants.
\end{proof}

\subsection{Local vertex mass and the temporal modulus}

\begin{proposition}[Uniform local mass bound]\label{prop:11-local-mass}
For every $\tau>0$ and $A>0$, there is $C_{\tau,A}<\infty$ such that
for all $H\ge2$, outside a set of $\mu_H$-mass at most
$C_{\tau,A}H^{-A}$, simultaneously for every visited vertex $v$ and
every real radius $1\le s\le H$,
\begin{equation}
 \#\bigl(V(\gamma)\cap\overline B(v,s)\bigr)\le H^\tau s^{4/3}.
 \label{eq:L-42}
\end{equation}
\end{proposition}

\begin{proof}
Use Theorem~\ref{thm:11-local-length} with a smaller slack $\tau'>0$.
It suffices to prove the statement at dyadic radii with further slack;
rounding a radius upwards then absorbs the constant for all radii.
There are polynomially many possible centers and $O(\log H)$ scales.
Small fixed radii are harmless by the lattice area bound.

For fixed $v,s$, surround $v$ by a lattice hexagonal box whose faces
are at perpendicular distances between $10s$ and $20s$. For each face
fix a parallel inner lattice line about $4s$ inward. Starting at the
next visit to $\overline B(v,s)$, follow the walk until its first subsequent hit
of a box face, or its endpoint; repeat at the next visit to the ball.
Each such visit batch is contained in the box. Its span is $O(s)$,
so the local length bound, with bounded trimming at face ports,
allows at most $CH^{\tau'}s^{4/3}$ visited vertices in that batch.
If \eqref{eq:L-42} fails, then for a fixed small
$0<\xi<\tau-\tau'$ and large $H$, at least
$a=\lceil H^\xi\rceil$ completed passages reach the same face.
There is at most one uncompleted terminal batch.

From each selected passage retain its segment from the last hit of
that face's inner line to the first hit of the face. These are
mutually disjoint strict bridges, all of height $j\asymp s$, in the
common box and hence a slab of width $Dj$, after a lattice symmetry.
If there are fewer than $a$ available ports, the event is empty.
Otherwise the endpoint tuples, matchings, and chronological placements have
at most $\exp(O(a\log H+\log H))$ choices. The complementary path
segments cost the same order: there are $a+1$ of them, each with
polynomial full-box susceptibility, and we may drop their mutual
avoidance. Apply Lemma~\ref{lem:11-many-bridges} to each fixed choice.
Since $s\le H$, the total mass is bounded by
\[
 \exp\!\left(Ca\log H+C\log H-\frac{ca^2}{\log H}\right),
\]
which is smaller than every power of $H^{-1}$ for $a=\lceil H^\xi\rceil$.
The polynomial union over centers and scales, and the exceptional
set for the local length estimate, preserve the claim.
\end{proof}

\begin{proposition}[Uniform temporal modulus]\label{prop:11-modulus}
For every $\tau>0$ and $A>0$ there is $C_{\tau,A}<\infty$ such that,
for every $H\ge2$, outside a set of $\mu_H$-mass at most
$C_{\tau,A}H^{-A}$, every subwalk $\sigma$ of $l$ steps satisfies
\[
 \operatorname{diam}(\sigma)\le H^\tau(1+l)^{3/4}.
\]
\end{proposition}

\begin{proof}
A violation has diameter $d>H^\tau(1+l)^{3/4}$. Some lattice normal
projection has span at least $cd$. Between its extreme levels take
an intervening weak bridge and extend it to a strict bridge. Its
height $h$ is at least $cd$ and at most $CH$, and its length is at
most $l+O(1)$. Thus $h\ge cH^\tau$ and, for a smaller fixed
$\tau'>0$ and large $H$,
\[
 l+O(1)\le C H^{-4\tau/3}h^{4/3}\le h^{4/3-\tau'}.
\]
Apply the fast-bridge estimate \eqref{eq:F-33b} dyadically in $h$.
The height lower bound is a positive power of $H$, and all placements
and exterior portions cost only a fixed polynomial by the full-box
bound. Its arbitrary failure power therefore gives the assertion.
\end{proof}
